\documentclass[../proyecto.tex]{subfiles}
\graphicspath{{\subfix{../images/}}}
\begin{document}

\chapter{Formatos}

% máximo 1 página

% Indicar cómo se espera abordar la tesis a nivel de escritura, y otras dimensiones extratextuales

Esta tesis incluye a grandes rasgos dos secciones fundamentales. La primera es el documento textual de la tesis doctoral que establece los fundamentos teóricos y conceptuales de la investigación realizada. La segunda es su ampliación con una colección de publicaciones disciplinares propias de la electrónica y la computación, que permiten el entendimiento de los procesos de diseño, fabricación y documentación de los artefactos técnicos de la tesis.

En la tesis se incluyen capítulos donde se describen los módulos necesarios para crear sintetizadores, incluyendo desde las metáforas, la historia y contexto cultural, hasta los diagramas de circuitos, diagramas de flujo, pseudocódigo y código fuente para su producción.

A nivel de acompañamiento electrónico, se incluyen diagramas de circuitos, listas de materiales, estrategias de diseño y fabricación de placas impresas, y los manuales técnicos de uso y reparación de sintetizadores y aparatos electrónicos.

A nivel de computación, se incluyen diagramas de flujo, pseudocódigo, y código fuente desarrollado con control de versiones. Este código es usado tanto para la síntesis de sonido dentro los sintetizadores realizados, como para la construcción de carcasas y otros artefactos físicos necesarios para esta empresa.

También se incluye la infraestructura empresarial y académica que ha permitido esta investigación, incluyendo la constitución de una sociedad por acciones para la compra de materiales y la venta de sintetizadores, los estudiantes pasantes o practicantes que han trabajado en este proyecto, y los cursos que han nutrido esta investigación, incluyendo los apuntes y libros escritos como material docente para estos cursos.

\end{document}
